\documentclass[11pt,a4paper]{report}

\usepackage{microtype}
\usepackage{xcolor}
\usepackage{subcaption}
\usepackage{amsmath}

% Editing macros — delete for the final draft
\newcommand{\fc}[1]{\textcolor{red}{[FactC: #1]}}
\newcommand{\ct}[1]{\textcolor{orange}{[CITE: #1]}}
\newcommand{\app}[1]{\textcolor{blue}{[APP: #1]}}
\newcommand{\sr}[1]{\textcolor{teal}{[Sreword: #1]}}
\newcommand{\fig}[1]{\textcolor{purple}{[FIG: #1]}}
\newcommand{\ai}[1]{\textcolor{brown}{[AI: #1]}}

\begin{document}
%=======================================================================================
\chapter{Data}
\section{Forest research data}
Data supplied by Forest Research as a GeoPackage.
Study area around Aberfoyle.
British National Grid, EPSG:27700.
Repeated forest-plot polygons from 2002, 2006, 2008, 2012, 2021 and 2023.
Each record represents a plot in one survey year.
Plots are roughly all 20m by 20m (area field median exactly 400\,m\textsuperscript{2} across all cohort plots), with smaller irregular exceptions (down to around 250\,m\textsuperscript{2}) where a plot's sampling grid cell is clipped by a compartment or sub-compartment boundary.
Plot and compartment are different levels of the same forestry hierarchy, not the same unit measured two ways: a compartment (\texttt{cpmt}) is a Forest Research management unit, further divided into sub-compartments (\texttt{scpt}); a plot is the individual LiDAR/Area-Based-Analysis sampling grid cell that survey statistics were computed over, and many plots sit inside one compartment.
The full data contain 795,645 plot-year records and 230,359 plot identifiers.
Sitka spruce was identified using \texttt{spis == "SS"}.
The main target was the 95th-percentile LiDAR height.
Important original variables included planting year, canopy cover, thinning information, species, yield class and windthrow hazard class.
The height data are Airborne Laser Scanning (ALS), collected from aircraft, not drone-based LiDAR.
%TODO: appendix: add full data dictionary or table in appendix

\begin{table}[h]
\centering
\begin{tabular}{p{2.2cm}p{3cm}p{3cm}p{3cm}p{3.5cm}}
\hline
Data family & Source & Resolution or structure & Temporal representation & Main purpose \\
\hline
Forest plots & Forest Research & Repeated plot polygons, roughly 20m by 20m & Six survey years, 2002--2023 & Prediction and growth calculation \\
Terrain & OS Terrain 50 & 50\,m grid & Static & Topography and exposure \\
Climate & HadUK-Grid & 1\,km grid (149 distinct values across 71,766 plots) & Annual or multi-year cohort-averaged summaries & Temperature, frost and rainfall conditions \\
Wind climate & Global Wind Atlas & Raster climatology (around 1,800 distinct values across the plots) & Long-term static summaries & General wind exposure \\
Observed wind & MIDAS stations & Station observations & Survey-interval summaries & Temporal wind and storm conditions \\
Other environment & CHELSA, CEH, SoilGrids and spatial layers & Source-dependent, 250\,m--1\,km & Mainly static & Wider environmental screening \\
\hline
\end{tabular}
\end{table}

\section{Data Cleaning}
The cleaning funnel builds two balanced cohorts from the same raw layer, applied in this order: complete coverage across every survey year of the chosen cohort, only Sitka spruce (\texttt{spis == "SS"}), a valid planting year (above 0, giving \texttt{Age = survey year - plyr}), then plausible age and height sanity checks.
In the current pipeline the age and height checks remove rows only from the four-survey cohort, and only at the planting-year step; the great majority of the attrition below comes from the completeness and species-selection stages.
The plausibility rules themselves: planting year must be no earlier than 1850 and strictly before the survey year; \texttt{Age} (survey year minus planting year) must fall between 0 and 200 years; height (the target column, \texttt{elev\_percentile\_95th}) must be greater than 0\,m and at most 60\,m, a generous ceiling for mature Sitka spruce.

\begin{table}[h]
\centering
\begin{tabular}{lrr}
\hline
Stage or characteristic & Four-survey cohort & Six-survey cohort \\
\hline
Complete plots before species filtering & 112,782 & 18,781 \\
Sitka spruce plots & 72,395 & 13,897 \\
Final plots & 71,766 & 13,897 \\
Final plot-year records & 287,064 & 83,382 \\
Survey years & 2008, 2012, 2021, 2023 & 2002, 2006, 2008, 2012, 2021, 2023 \\
Observations per plot & 4 & 6 \\
Survey intervals & 4, 9 and 2 years & 4, 2, 4, 9 and 2 years \\
Role & Main analysis & Sensitivity analysis \\
\hline
\end{tabular}
\end{table}

The four-survey cohort has more plots, so greater spatial coverage.
The six-survey cohort has two more observations per plot but retained fewer than one-fifth as many plots.
The four-survey cohort was therefore used for the main analysis; the six-survey cohort tested whether more years changed the conclusions.

\section{stand terrain and environmental data}
\subsection{stand and management}
Canopy cover represented the observed canopy condition.
Recorded thinning information was used to derive thinning status, time since thinning and recent-thinning indicators.
Yield class was inspected but excluded from the main prediction features because it is derived from forest growth information and presents a leakage or circularity concern.
Windthrow hazard class was treated as an existing forest descriptor rather than a direct observation of wind exposure, and was not included in the models until wind data were also added, since it showed a similar spatial relationship and was constant across survey years for a given plot.

These stand-management variables entered the main stand-based DNN and PINN comparison as the no-environment feature set.
Most environmental and terrain data (elevation, wind, CHELSA climatologies, soil) entered both the prediction models (\texttt{dnn\_env\_terrain}, \texttt{pinn\_env\_terrain}) and the later attribution analyses, so prediction and attribution were not fully separate questions for those variables.
HadUK-Grid climate variables (\texttt{tas\_mean}, \texttt{groundfrost\_mean}) are the one exception: they are currently used only in the attribution-side screening (XGBoost/SHAP), because of a structural gap in the prediction pipeline's handling of cohort-specific column names, not because the variables were judged untrustworthy.

\section{Terrain data}
OS Terrain 50 provided a static representation of topography at approximately 50\,m resolution.

The derived variables can be grouped by meaning: height and steepness (elevation and slope); orientation (aspect, and the northness/eastness measures derived from it); landform (terrain position index, curvature and local relief); exposure and shelter (TOPEX and related horizon measures); and local terrain heuristics (the frost-hollow flag and any investigated soil-depth proxy).

The raw aspect field is a 0--360\textdegree{} compass bearing and was not used directly as a model input, since a circular bearing has a discontinuity at 0\textdegree{}/360\textdegree{} that a model would otherwise treat as a large numeric jump between two directions that are physically adjacent.
It was instead replaced by its two continuous trigonometric components, computed from the aspect in radians ($\theta$):
\begin{align*}
\text{northness} &= \cos(\theta) \\
\text{eastness} &= \sin(\theta)
\end{align*}
so a north-facing slope (aspect $0\textdegree$) has northness $=1$, and an east-facing slope (aspect $90\textdegree$) has eastness $=1$.

Terrain is static and cannot explain changes in conditions between surveys.
TOPEX describes topographic exposure or shelter; it is not measured wind.
The frost-hollow flag is a project-derived indicator, not an observation of frost damage.
Any proposed soil-depth proxy should be labelled as exploratory unless it has strong ecological support.

\section{Wind data}
Global Wind Atlas provides raster-based, long-term wind-climate summaries.
It represents broad spatial differences in typical wind conditions and is effectively static for this analysis.
It can support spatial attribution but cannot identify conditions during a specific survey interval.

MIDAS station observations provide temporally varying weather-station measurements.
These were summarised over the intervals between LiDAR surveys, allowing investigation of mean wind, extreme wind or storm conditions during different growth periods.
They may not represent conditions at every forest plot equally, since station observations require spatial matching or interpolation.

Global Wind Atlas variables represented long-term spatial wind climate, MIDAS variables represented observed conditions over particular periods, and TOPEX represented terrain-based exposure.
These sources describe related but different aspects of wind exposure and were not treated as interchangeable.

\section{Climate, soil and other spatial data}
HadUK-Grid provided data at approximately 1\,km spatial resolution, used to investigate variables such as temperature, ground frost and rainfall.
Depending on the extraction, variables represented individual years, or were averaged across each cohort's own survey years once the multi-year extraction was in place.
The 1\,km resolution means many nearby plots share identical values.
As a compact diagnostic example: the temperature variable contained only 149 distinct values across 71,766 plots, and showed strong spatial structure, including Moran's I $=0.941$ and a between-compartment intraclass correlation of 0.893.
It therefore represented broad spatial climate zones more strongly than fine-scale variation between neighbouring plots.
%METHODOLOGY, not here: the Spearman-vs-residual screening number and the reasoning for why this variable was kept despite its coarseness belong in the feature-selection part of Methodology (ahead of 4.4), alongside the same before/after Moran's I check applied to model residuals for every other candidate feature.

ERA5-Land was investigated as an additional reanalysis source.
Its temperature extraction produced only eight distinct values across the study plots, coarser than HadUK-Grid's 149.
It was not retained in the final feature sets used for modelling.

CHELSA supplied long-term (1981--2010) climatic summaries at the same nominal 1\,km grid as HadUK-Grid, but using terrain-conditioned downscaling rather than plain interpolation.
It was retained in the final feature sets used for both the attribution screening and the environment/terrain prediction models.

SoilGrids pH and the CEH soil and terrain classifications were explored as possible controls on spatial growth variation.
These are largely static variables that may describe broad site conditions but cannot establish short-term environmental effects.
The CEH categorical layers (pedotope, subsurface drainage, textural composition) each carry only 3 to 7 classes across the plots, so the same class is assigned to many plots.

Distances to roads, watercourses and forest boundaries were explored where relevant.
These are spatial context variables, not direct measures of growth processes.

\begin{table}[h]
\centering
\begin{tabular}{p{5.5cm}p{7.5cm}}
\hline
Limitation & Effect on interpretation \\
\hline
Complete-observation requirement & May favour stable plots that remained identifiable across all surveys \\
Four or six observations per plot & Limits estimation of detailed plot-specific trajectories \\
Unequal survey intervals & Requires annualisation and may hide within-interval changes \\
LiDAR height is not identical to biological growth & Changes may include canopy, measurement or disturbance effects \\
Limited acquisition metadata & Survey-year measurement differences cannot be fully separated \\
Coarse environmental rasters & Many plots receive the same climate value (e.g.\ 149 distinct HadUK-Grid values, 8 distinct ERA5-Land values, across 71,766 plots) \\
Static terrain and climate summaries & Describe spatial conditions more readily than temporal change \\
Station-based wind observations & Conditions at a station may differ from conditions at individual plots \\
Uncertain management history & Recorded thinning may not capture every intervention or disturbance \\
Multiple correlated environmental variables & Makes individual attribution difficult and requires controlled selection \\
Observational data & Environmental relationships indicate association, not causation \\
\hline
\end{tabular}
\end{table}

The environmental datasets provide several representations of climate, terrain and exposure, but their resolutions and time periods differ.
The attribution analysis can therefore identify variables associated with departures from expected growth, but it cannot isolate causal environmental effects.

\end{document}
